\documentclass[UTF8,aspectratio=169]{ctexbeamer}
\usetheme{Madrid}
\usecolortheme{default}
\usepackage{amsmath}
\usepackage{booktabs}
\usepackage{siunitx}
\usepackage{graphicx}
\usepackage{array}
\usepackage{xcolor}
\usepackage{xurl}

\title[Stage-3 Raw CDF Refinement]{Stage-3 Raw CDF Refinement\\Distillation, Ablation, and Final Metrics}
\author{Mie Spray Postprocessing Project}
\date{\today}

\newcommand{\good}[1]{\textcolor{green!45!black}{#1}}
\newcommand{\warn}[1]{\textcolor{red!65!black}{#1}}
\newcommand{\src}[1]{\vfill{\tiny\textcolor{gray!75!black}{Source: \url{#1}}}}
\newcolumntype{L}[1]{>{\raggedright\arraybackslash}p{#1}}
\newcolumntype{R}[1]{>{\raggedleft\arraybackslash}p{#1}}

\begin{document}

\begin{frame}
  \titlepage
\end{frame}

\begin{frame}{Message}
\small
\textbf{Claim:} Stage-3 refinement turns the Stage-2 heteroscedastic teacher into a raw-CDF-aware predictor by mixing direct raw supervision with teacher distillation under time-dependent coverage regimes.

\vspace{0.45em}
\textbf{Current status after the \(\Delta P^{0.5}\) feature-engineering update:}
\begin{itemize}
  \item the teacher/student feature space uses \(A_{\mathrm{scale}}=\Delta P^{0.5}\rho_a^{-0.25}d_n^{0.5}\);
  \item production Stage-3 is evaluated over five seeds: 7/17/42/99/2024;
  \item the dominant winning ablation is \texttt{anchor\_off};
  \item refreshed full-clean evaluation has \num{692942} points and \num{24316} trajectories;
  \item empirical baselines and the Stage~3 SVGP are included for comparison.
\end{itemize}
\end{frame}

\begin{frame}{Why Stage-3 Exists}
\small
\begin{block}{Stage-2 teacher}
Stage 2 learns a heteroscedastic Gaussian model from fitted representative curves. It is smooth and defined everywhere on \(0\)--\SI{5}{ms}, but it can inherit fitting assumptions from the reduced curve model.
\end{block}

\begin{block}{Raw CDF series}
Raw CDF observations are closer to measured data, but they are incomplete and time-dependent: coverage thins as trajectories disappear from the field of view or fail quality filters.
\end{block}

\begin{block}{Stage-3 refinement}
Stage 3 uses raw CDF where it is reliable, teacher distillation where raw data are sparse, and light shape/onset priors to keep trajectories physically plausible.
\end{block}
\end{frame}

\begin{frame}{CDF Coverage Regimes}
\small
Time is binned into \(\Delta t=\SI{0.1}{ms}\) intervals. For each operating-condition group, smoothed CDF sample counts define three supervision regimes:

\vspace{0.4em}
\begin{tabular}{L{0.24\linewidth}L{0.28\linewidth}L{0.36\linewidth}}
\toprule
Regime & Trigger & Supervision role \\
\midrule
\texttt{raw\_reliable} & coverage \(\geq70\%\) of peak & raw CDF is trusted directly \\
\texttt{raw\_uncertain} & below \(70\%\), above \(20\%\) & teacher dominates transition \\
\texttt{teacher\_only} & below \(20\%\) or \(<4\) samples & teacher is the only usable signal \\
\bottomrule
\end{tabular}

\vspace{0.5em}
Transitions require consecutive bins, which prevents isolated coverage dips from changing the regime.
\end{frame}

\begin{frame}{Default Supervision Weights}
\small
\begin{center}
\begin{tabular}{L{0.32\linewidth}R{0.20\linewidth}R{0.20\linewidth}}
\toprule
Regime & Raw NLL weight & KD KL weight \\
\midrule
\texttt{raw\_reliable} & 1.00 & 0.25 \\
\texttt{raw\_uncertain} & 0.00 & 1.00 \\
\texttt{teacher\_only} & 0.00 & 0.75 \\
\bottomrule
\end{tabular}
\end{center}

\vspace{0.45em}
\begin{itemize}
  \item Raw NLL is used only where raw CDF is reliable and finite.
  \item KD KL prevents the student from drifting where raw observations are sparse.
  \item A global inverse-frequency time-bin weight is clipped to \(0.5\)--\(2.0\).
  \item Teacher confidence further down-weights high-\(\sigma\) teacher regions.
\end{itemize}

\src{MLP/MLP_training/train_stage3_distillation_plus_raw_series.py}
\end{frame}

\begin{frame}{\(\Delta P^{0.5}\) A-Scaled Physical Space}
\small
\textbf{Current engineered scale:}
\[
A_{\mathrm{scale}}=\Delta P^{0.5}\rho_a^{-0.25}d_n^{0.5}.
\]

\textbf{Network output is still scaled:}
\[
\hat \mu(t)=\mu(t)/A,\qquad \hat\sigma(t)=\sigma(t)/A.
\]

\textbf{Stage-3 losses are computed in physical millimetres:}
\[
\mu_{\mathrm{phys}}=A\hat\mu,\qquad
\sigma_{\mathrm{phys}}=A\hat\sigma.
\]

\vspace{0.25em}
This lets raw CDF targets enter directly in mm while keeping the sparse pressure/diameter axes collapsed through the data-supported \(A\)-scale.

\src{MLP/MLP_training/engineered_feature_common.py}
\end{frame}

\begin{frame}{Composite Loss}
\small
\[
\mathcal{L}
= \mathcal{L}_{\mathrm{raw}}
+ \mathcal{L}_{\mathrm{KD}}
+ \lambda_{\mathrm{onset}}\mathcal{L}_{\mathrm{onset}}
+ \lambda_{\mathrm{anchor}}\mathcal{L}_{\mathrm{anchor}}
+ \lambda_{d_1}P_{d_1}
+ \lambda_{d_2}P_{d_2}.
\]

\vspace{0.2em}
\begin{tabular}{L{0.23\linewidth}L{0.68\linewidth}}
\toprule
Term & Role \\
\midrule
Raw NLL & Gaussian NLL against raw CDF in physical space \\
KD KL & Gaussian KL between student and Stage-2 teacher in physical space \\
Onset BCE & soft ramp target in the first \(\SI{0.2}{ms}\) \\
Anchor & optional early-time penalty on \(\mu(t\approx0)\) \\
\(P_{d_1}\) & monotonicity penalty against negative increments \\
\(P_{d_2}\) & concavity penalty after \(t\geq\SI{0.5}{ms}\) \\
\bottomrule
\end{tabular}
\end{frame}

\begin{frame}{Loss Terms: Probabilistic Meaning}
\small
\begin{block}{Raw Gaussian NLL}
\footnotesize
Where raw CDF samples are reliable, the student is scored directly in physical mm:
\[
\mathcal{L}_{\mathrm{raw},i}
=\frac{1}{2}\left[
\frac{(y_i-\mu_i)^2}{\sigma_i^2}
+\log\sigma_i^2
\right] + \mathrm{const.}
\]
This rewards both accurate means and correctly sized variances; shrinking
\(\sigma_i\) is penalized if residuals remain large.
\end{block}

\begin{block}{Teacher--student Gaussian KL}
\footnotesize
When raw coverage is weak, the student is kept close to the Stage-2 teacher:
\[
D_{\mathrm{KL}}(T\Vert S)
=\log\frac{\sigma_S}{\sigma_T}
+\frac{\sigma_T^2+(\mu_T-\mu_S)^2}{2\sigma_S^2}
-\frac{1}{2}.
\]
The direction \(T\Vert S\) treats the teacher distribution as the reference
distribution that the student should cover in sparse regions.
\end{block}
\end{frame}

\begin{frame}{Key Ablation Finding: Anchor Off}
\small
\begin{block}{Old design assumption}
The early anchor was originally included to keep \(\mu(t\approx0)\) close to zero and prevent unphysical onset behavior.
\end{block}

\begin{block}{What the ablation found}
The best current Stage-3 setting is usually \texttt{anchor\_off}, i.e. \(\lambda_{\mathrm{anchor}}=0\). In the five-seed \(\Delta P^{0.5}\) mode-C pipeline, seeds 42, 17, 99, and 7 selected \texttt{anchor\_off}; seed 2024 selected \texttt{raw\_reliable\_no\_kd}.
\end{block}

\begin{block}{Interpretation}
The learned onset head, raw/KD regime weights, and shape priors are enough to regularize early time. A hard \(\mu(0)\) anchor can overconstrain the raw-CDF refinement.
\end{block}
\end{frame}

\begin{frame}{Metric Contract for Stage-3 Results}
\small
\begin{block}{Accuracy metrics}
\footnotesize
RMSE emphasizes large misses; MAE reports the typical absolute millimetre error;
P95 absolute error reports the tail. Bias keeps the sign, so it exposes systematic
over- or under-prediction that RMSE/MAE hide.
\end{block}

\begin{block}{Uncertainty metrics}
\footnotesize
Coverage is
\[
\Pr(|y-\mu|\leq k\sigma),\qquad k\in\{1,2\}.
\]
For a calibrated Gaussian, the nominal targets are 0.683 and 0.954. Higher values
mean conservative intervals; lower values mean under-coverage.
\end{block}

\begin{block}{Protocol separation}
\footnotesize
The five-seed in-pipeline result measures training stability on each seed's
post-evaluation subset. The external clean diagnostic is the headline comparison
because every model is evaluated on the same \num{692942} full-clean points or
the same \num{542565} conservatively uncensored CDF points.
\end{block}
\end{frame}

\begin{frame}{Production Five-Seed Repeatability}
\small
This is the refreshed production \texttt{anchor\_off} Stage-3 model evaluated
over seeds 7/17/42/99/2024.

\vspace{0.35em}
\begin{center}
\begin{tabular}{L{0.28\linewidth}R{0.24\linewidth}R{0.24\linewidth}}
\toprule
Metric & Full clean & CDF uncensored \\
\midrule
RMSE mm & 4.735 \(\pm\) 0.037 & 4.265 \(\pm\) 0.040 \\
MAE mm & 3.293 \(\pm\) 0.039 & 3.032 \(\pm\) 0.048 \\
Bias mm & +0.626 \(\pm\) 0.112 & +0.562 \(\pm\) 0.119 \\
P95 abs. err. mm & 9.167 \(\pm\) 0.066 & 8.307 \(\pm\) 0.076 \\
1\(\sigma\) coverage & 0.867 \(\pm\) 0.004 & 0.887 \(\pm\) 0.004 \\
2\(\sigma\) coverage & 0.985 \(\pm\) 0.001 & 0.989 \(\pm\) 0.001 \\
\bottomrule
\end{tabular}
\end{center}

\vspace{0.3em}
\footnotesize
The per-seed spread is small; the conclusion is not seed luck. The CDF-uncensored
table is the cleaner physical metric, while full-clean preserves continuity with
earlier thesis material.

\src{Thesis/generated/baseline_comparison_20260521/fixed_table_group_summary.csv}
\end{frame}

\begin{frame}{External Clean Diagnostic}
\small
The refreshed production Stage-3 models were evaluated on the full clean diagnostic protocol:
\num{692942} points and \num{24316} trajectories.

\vspace{0.35em}
\begin{center}
\begin{tabular}{L{0.38\linewidth}R{0.12\linewidth}R{0.12\linewidth}R{0.12\linewidth}R{0.12\linewidth}}
\toprule
Model & RMSE & MAE & P95 & 1\(\sigma\) cov. \\
\midrule
Hiroyasu--Arai calibrated & 11.268 & 8.658 & 22.891 & 0.615 \\
Naber--Siebers delay & 10.545 & 8.177 & 21.253 & 0.637 \\
\good{Production MLP anchor-off, 5-seed mean} & \good{4.735} & 3.293 & \good{9.167} & \good{0.867} \\
\(\mu\)-anchor raw-uncertain MLP, 5-seed mean & 5.200 & 3.654 & 10.525 & 0.834 \\
Sparse heteroscedastic GP, Stage~3 & \good{4.628} & \good{3.173} & 9.242 & 0.703 \\
\bottomrule
\end{tabular}
\end{center}

\vspace{0.25em}
The production MLP crushes the physical baselines and is close to SVGP, but the
SVGP remains the point-accuracy winner. The MLP contributes conservative
coverage and a deployable neural surrogate.

\src{Thesis/generated/baseline_comparison_20260521/headline_full_clean.csv}
\end{frame}

\begin{frame}{Clean vs Hard-Subset Calibration}
\small
\begin{block}{External clean diagnostic}
The production Stage-3 MLP over-covers on the clean full diagnostic:
1\(\sigma\)/2\(\sigma\) coverage is \(0.867/0.985\), above nominal \(0.683/0.954\).
\end{block}

\begin{block}{In-pipeline post-evaluation subset}
The same pipeline has 1\(\sigma\)/2\(\sigma\) coverage of \(0.654/0.966\), much closer to nominal and slightly under at 1\(\sigma\).
\end{block}

\begin{block}{Paper interpretation}
The learned \(\sigma\) is not a single global error bar. It behaves differently on clean/easy trajectories and harder representative rows, which is evidence for difficulty-aware uncertainty rather than pure input-independent noise.
\end{block}
\end{frame}

\begin{frame}{Raw Supervision Source Ablation (single run)}
\small
Does replacing \texttt{series\_wide\_clean} with the per-condition
right-censored \texttt{series\_wide\_truncated} improve generalisation?
\vspace{0.35em}
\begin{center}
\begin{tabular}{L{0.38\linewidth}R{0.13\linewidth}R{0.13\linewidth}R{0.13\linewidth}R{0.13\linewidth}}
\toprule
Configuration & RMSE & MAE & Bias & P95 \\
\midrule
Production (\texttt{series\_wide\_clean})     & 4.746 & 3.290 & --- & --- \\
Truncated source (\texttt{series\_wide\_truncated}) & 4.743 & 3.305 & +0.751 & 9.228 \\
\midrule
\(\Delta\) (truncated \(-\) production) & \(-\)0.003 & +0.015 & --- & --- \\
\bottomrule
\end{tabular}
\end{center}
\vspace{0.3em}
\textbf{Conclusion.} The \(\Delta\)RMSE is negligible (\(\approx0.003\) mm).
The regime weight mechanism already zeros the raw loss for low-confidence
bins; explicitly truncating those frames in the source table is therefore
redundant. This serves as a positive control confirming that the regime
labelling correctly isolates the unreliable frames.

\src{MLP/MLP\_training/stage3\_truncated\_supervision\_config.json}
\end{frame}

\begin{frame}{Regime Weight Filter Ablation (single run)}
\small
What is the cost of \emph{not} zeroing \texttt{raw\_weight} for the
\texttt{raw\_uncertain} and \texttt{teacher\_only} regimes?
\vspace{0.35em}
\begin{center}
\begin{tabular}{L{0.38\linewidth}R{0.13\linewidth}R{0.13\linewidth}R{0.13\linewidth}R{0.13\linewidth}}
\toprule
Configuration & RMSE & MAE & Bias & P95 \\
\midrule
Production (filter \textbf{on}: weights = 0) & 4.746 & 3.290 & --- & --- \\
No regime filter (weights = 1.0) & 4.811 & 3.350 & +0.721 & 9.254 \\
\midrule
\(\Delta\) (no-filter \(-\) production) & +0.065 & +0.060 & --- & +0.026 \\
\bottomrule
\end{tabular}
\end{center}
\vspace{0.3em}
\textbf{Conclusion.} Admitting right-censored late-time frames as raw
supervision degrades RMSE by \(+0.065\) mm (\(\approx1.4\%\)).
The KD term partially compensates (high KD weights in uncertain bins
suppress the noise), limiting the damage but not eliminating it.
Zeroing the raw weight in low-confidence bins is therefore a justified
and quantifiably beneficial design choice.

\src{MLP/MLP\_training/stage3\_no\_regime\_raw\_filter\_config.json}
\end{frame}

\begin{frame}{What Should Go Into the Thesis}
\small
\begin{block}{Method sentence}
Stage 3 refines the Stage-2 heteroscedastic teacher against raw CDF measurements using regime-dependent raw NLL and teacher KL losses, with regimes determined by per-condition CDF coverage over time.
\end{block}

\begin{block}{Ablation sentence}
The best \(\Delta P^{0.5}\) Stage-3 configuration removes the explicit early-time anchor in four of five seeds, indicating that the onset head and raw/KD regime weighting provide sufficient early-time regularization. Two additional single-run ablations confirm that the regime filter is effective: replacing the clean source with the right-censored truncated source yields negligible change (\(\Delta\)RMSE \(\approx-0.003\) mm), while disabling the regime raw-weight filter degrades RMSE by \(+0.065\) mm.
\end{block}

\begin{block}{Result sentence}
On the refreshed full-clean benchmark, the resulting five-seed production MLP
reaches RMSE \(4.735\pm0.037\) mm and outperforms Hiroyasu--Arai and
Naber--Siebers by more than 55\% in RMSE. It is slightly behind the Stage~3
SVGP in point RMSE, but gives more conservative coverage.
\end{block}
\end{frame}

\begin{frame}{Data Sources for Thesis Extraction}
\small
\begin{tabular}{L{0.30\linewidth}L{0.60\linewidth}}
\toprule
Item & Path \\
\midrule
Stage-3 training script & \scriptsize\url{MLP/MLP_training/train_stage3_distillation_plus_raw_series.py} \\
Production MLP runs & \scriptsize\url{MLP/runs_mlp/distill_cdf_onset_v2_ablate_anchor_off_20260521_122435}--\url{..._124004} \\
Fixed-table comparison & \scriptsize\url{MLP/baseline/comparison_reports/stage3_fixed_table_eval_with_ha_ns_20260521/} \\
Full-clean comparison table & \scriptsize\url{Thesis/generated/baseline_comparison_20260521/headline_full_clean.csv} \\
SVGP Stage~3 output & \scriptsize\url{MLP/runs_mlp/gp_baseline_stage3_20260521_112229/} \\
\bottomrule
\end{tabular}
\end{frame}

\end{document}
